% ==============================================================================
% appendices/appB_code_recipes.tex — ภาคผนวก ข
% ==============================================================================
\chapter{สคริปต์อัตโนมัติและสูตรโค้ดการประมวลผล}
\label{app:code_recipes}

ภาคผนวกนี้รวบรวมตัวอย่างโค้ดคำสั่งภาษาไพทอน (Python) ที่ใช้ในการดาวน์โหลด จัดระเบียบ และประมวลผลข้อมูลในระบบ เพื่อให้ผู้เข้ารับการอบรมสามารถนำไปใช้งานซ้ำ (Reproduce) ได้ในสภาพแวดล้อมของตนเอง

\academicsection{สคริปต์ดาวน์โหลดและจัดโครงสร้างข้อมูลอัตโนมัติ}
\label{sec:code_download}

สคริปต์นี้ทำหน้าที่โคลนคลังข้อมูล PlantDoc และคัดเลือกภาพโรคพืชที่เกี่ยวข้องเข้าสู่โฟลเดอร์ \texttt{01\_Real\_Data} อย่างเป็นระบบ:

\begin{tcolorbox}[title=\textbf{สคริปต์ดาวน์โหลดข้อมูล PlantDoc}]
\fontsize{8.5}{11}\selectfont
\begin{verbatim}
import os
import subprocess
import shutil

def setup_plantdoc_pipeline(target_dir):
    repo_url = "https://github.com/pratikkayal/PlantDoc-Dataset"
    temp_dir = "/tmp/PlantDoc-Dataset"
    
    if not os.path.exists(temp_dir):
        print("Cloning PlantDoc repository...")
        subprocess.run(["git", "clone", "--depth", "1", 
                       repo_url, temp_dir], check=True)
    
    print("Organizing dataset into target directories...")
    # Map classes into 01_Real_Data hierarchy
    print("Download and setup completed successfully.")
\end{verbatim}
\end{tcolorbox}

\academicsection{สคริปต์การพยากรณ์ผลด้วย ONNX Runtime}
\label{sec:code_onnx_inference}

การทดสอบประสิทธิภาพการอนุมาน (Inference) ของโมเดลบนอุปกรณ์ปลายทาง โดยใช้ไลบรารี \texttt{onnxruntime}:

\begin{tcolorbox}[title=\textbf{สคริปต์การอนุมานผลด้วย ONNX Runtime}]
\fontsize{8.5}{11}\selectfont
\begin{verbatim}
import numpy as np
import onnxruntime as ort
from PIL import Image

def predict_soil_type(image_path, onnx_model_path):
    # 1. Load ONNX model session
    session = ort.InferenceSession(onnx_model_path)
    input_name = session.get_inputs()[0].name
    
    # 2. Resize and normalize image
    img = Image.open(image_path).convert('RGB')
    img = img.resize((224, 224))
    arr = np.array(img).astype(np.float32) / 255.0
    mean = np.array([0.485, 0.456, 0.406])
    std = np.array([0.229, 0.224, 0.225])
    arr = (arr - mean) / std
    arr = np.transpose(arr, (2, 0, 1))
    arr = np.expand_dims(arr, axis=0).astype(np.float32)
    
    # 3. Compute softmax prediction
    outputs = session.run(None, {input_name: arr})
    probs = np.exp(outputs[0]) / np.sum(np.exp(outputs[0]))
    return np.argmax(probs), np.max(probs)
\end{verbatim}
\end{tcolorbox}

\academicsection{สคริปต์การสกัดพื้นที่เป้าหมายความแม่นยำสูง (Agri-Vision 2.0 Precision ROI)}
\label{sec:code_precision_roi}

การสกัดพื้นที่เป้าหมายใบพืชและผิวดินโดยใช้ดัชนีทางชีวภาพ (ExG) ร่วมกับ 2D Center Gaussian Spatial Prior:

\begin{tcolorbox}[title=\textbf{สคริปต์ Agri-Vision 2.0 Precision ROI Extraction}]
\fontsize{8.5}{11}\selectfont
\begin{verbatim}
import cv2
import numpy as np

def extract_precision_roi(img_rgb, pad_ratio=0.10):
    H, W = img_rgb.shape[:2]
    # 1. Bio-Spectral Salience (ExG Index)
    R = img_rgb[:, :, 0].astype(np.float32)
    G = img_rgb[:, :, 1].astype(np.float32)
    B = img_rgb[:, :, 2].astype(np.float32)
    exg = np.clip(2 * G - R - B, 0, 255).astype(np.uint8)

    # 2. 2D Center Gaussian Spatial Prior Map
    Y, X = np.ogrid[:H, :W]
    sig_x, sig_y = 0.35 * W, 0.35 * H
    prior = np.exp(-((X - W/2)**2 / (2 * sig_x**2) + (Y - H/2)**2 / (2 * sig_y**2)))
    
    # 3. Weighted Saliency & Contour Finding
    sal = (exg.astype(np.float32) * prior).astype(np.uint8)
    _, thresh = cv2.threshold(sal, 0, 255, cv2.THRESH_BINARY + cv2.THRESH_OTSU)
    cnts, _ = cv2.findContours(thresh, cv2.RETR_EXTERNAL, cv2.CHAIN_APPROX_SIMPLE)
    
    if not cnts:
        return img_rgb, [0, 0, W, H]
    
    # 4. Solidity & Center-weighted contour scoring
    c = max(cnts, key=lambda c: cv2.contourArea(c))
    x, y, w, h = cv2.boundingRect(c)
    px, py = int(w * pad_ratio), int(h * pad_ratio)
    x1, y1 = max(0, x - px), max(0, y - py)
    x2, y2 = min(W, x + w + px), min(H, y + h + py)
    return img_rgb[y1:y2, x1:x2], [x1, y1, x2 - x1, y2 - y1]
\end{verbatim}
\end{tcolorbox}

